%% ma: L11-B15-0002
%% lop: 11
%% chuong: 5
%% bai: 15
%% dang: 11.15.1
%% loai: TL
%% muc: TH
%% dap_an: "-"
%% nguon: "tai-lieu-hoc-tap-mon-toan-lop-11-theo-chuong-trinh-giao-duc-pho-thong-2018.pdf (Lớp Toán Thầy Duy, Phạm Lê Duy, THPT Chu Văn An), Chương 3 Bài 1, Dạng 1, Ví dụ 2 (tr.307)"
%% kien_thuc: "Lớp 11 Bài 15 (giới hạn của dãy số)"
%% trang_thai: cho-duyet
%% ly_do: "Dạng mới đề xuất (11.15.x), chờ giáo viên duyệt dạng; tài liệu gốc không có lời giải mẫu, đáp án do kho tự giải và kiểm bằng sympy."
\begin{ex}
Chứng minh rằng

a) $\lim\dfrac{2n^2+n}{n^2+4}=2$;

b) $\lim\dfrac{6n+2}{n+5}=6$;

c) $\lim\dfrac{7^n-2\cdot8^n}{8^n+3^n}=-2$;

d) $\lim\dfrac{2\cdot3^n+5^n}{5^n+3^n}=1$.
\loigiai{\sloppy a) $\dfrac{2n^2+n}{n^2+4}-2=\dfrac{n-8}{n^2+4}=\dfrac{\frac1n-\frac8{n^2}}{1+\frac4{n^2}}\to0$.

b) $\dfrac{6n+2}{n+5}-6=\dfrac{-28}{n+5}\to0$.

c) \[\begin{aligned}\dfrac{7^n-2\cdot8^n}{8^n+3^n}+2=\dfrac{7^n+2\cdot3^n}{8^n+3^n}\\ &=\dfrac{\left(\frac78\right)^n+2\left(\frac38\right)^n}{1+\left(\frac38\right)^n}\to0.\end{aligned}\]

d) \[\dfrac{2\cdot3^n+5^n}{5^n+3^n}-1=\dfrac{3^n}{5^n+3^n}=\dfrac{\left(\frac35\right)^n}{1+\left(\frac35\right)^n}\to0.\]}
\end{ex}
